\setupcolors[state=start]
\setupcolor[hex]
\definecolor[darkblue][h=003366]
\definecolor[darkgray][h=e6e6e6]
\definecolor[redline][h=d9534f]
\definecolor[blueline][h=5bc0de]
\definecolor[greenline][h=5cb85c]

\setuphead[section][style=\bfa, color=darkblue, before=\blank, after=\blank]
\setuphead[subsection][style=\bf, color=black, before=\blank, after=\blank]
\setuppapersize[A4][A4]
\setupbodyfont[11pt]
\setuplayout[width=middle, height=middle, topspace=2cm, bottomspace=2cm]
\mainlanguage[es]

\starttext
\title{\bf Exploración de OPP\_v5\_2\_split: qué encarece cada etapa y dónde deja de resolver}
\blank[big]

Este documento reúne los siete estudios exploratorios que se corrieron sobre el modelo
OPP\_v5\_2\_split el 22 y el 23 de septiembre de 2026. La pregunta era práctica: los tests de
stress heredados de OPP\_v4 no terminaban nunca, y hacía falta saber qué dimensión de la
instancia vuelve difícil a cada etapa del modelo para poder diseñar escenarios de stress que
fueran realistas y a la vez resolubles.

\blank

\startsection[title={Alcance y fuentes}]
\startitemize
\item {\bf Modelo:} la exploración es de OPP\_v5\_2\_split. OPP\_v5\_1\_mono aparece en dos
puntos solamente, cada uno comparado contra el split sobre la misma instancia
(sección \in[sec:mono]).
\item {\bf Datos:} un directorio por estudio en \type{experiments/opp_v5/studies/}, con los
puntos o scripts que se corrieron y los CSV que produjeron. La sección \in[sec:tests] suma las
corridas de los tests de stress del perfil realista (\type{results/opp_v5_2_split/stress/}),
que no son exploración pero validan sus conclusiones.
\item {\bf Instancias:} todas salen del mismo generador (\type{designed_stress_scenario}):
grilla de depósito con pesos vertical/horizontal, ubicaciones de las órdenes al azar con
semilla fija, lotes armados por turnos, fechas de entrega repartidas entre el 60\% y el 140\%
de una estimación del makespan, y capacidad del vehículo holgada para que nunca sea la
restricción activa.
\item {\bf No son tests:} un test responde si una instancia resuelve; un barrido da la curva,
es decir dónde se dobla cada eje y dónde se rompe.
\stopitemize
\stopsection

\startsection[title={Cómo leer los tiempos}]
Todo se midió en una notebook con un Intel i7-1360P: 4 núcleos de rendimiento con
hyperthreading y 8 de eficiencia, que van a cerca de la mitad de velocidad. Eso condiciona cómo
se leen los números.

\startitemize
\item {\bf Los nodos de branch-and-bound son deterministas.} SCIP recorre el mismo árbol ante el
mismo input. La instancia de 6 lotes, 3 pickers y 26 órdenes por lote dio 14.041 nodos en dos
corridas distintas, con 12,9 s y 17,7 s de reloj. Con 4 pickers, 31.317 nodos en las dos, con
25,4 s y 26,3 s. Para comparar dificultad entre corridas hechas en condiciones distintas, los
nodos son la medida justa.
\item {\bf El reloj varía hasta 2 veces.} La calibración del perfil realista, medida de a un punto,
tardó 214 s en 7×7 y 710 s en 13×13; los tests de stress sobre esas mismas instancias, un rato
después, 123 s y 351 s, con los mismos objetivos.
\item {\bf Correr en paralelo infla el reloj entre 1,6 y 2,3 veces} (500 s contra 214 s y
1.128 s contra 710 s, con los mismos nodos). Solo las comparaciones lado a lado dentro de una
misma corrida son confiables en tiempo.
\item {\bf El límite de tiempo de SCIP no es un tope duro.} SCIP lo consulta entre operaciones.
Un punto con \type{limits/time = 300} terminó a los 3.079 s, con 499 nodos, porque cada LP
tardaba varios segundos. Por eso cada punto corre en su propio proceso y se mata desde afuera
(estudio 04).
\item {\bf Los estudios 01 a 03 corrieron con la clausura métrica activada} (sección
\in[sec:clausura]). La clausura es exacta: no cambia ni el objetivo ni las rutas, solo el tiempo
de la Etapa 1. Los tiempos de Etapa 1 de esos estudios son comparables entre sí, pero no
directamente con los de los estudios posteriores.
\stopitemize
\stopsection

\startsection[title={El modelo en dos etapas}]
La Etapa 1 resuelve, para cada lote por separado, el camino de menor costo que sale del depósito,
visita todas las ubicaciones del lote y termina en el nodo final. La conectividad la garantiza un
flujo de Gavish-Graves ($q_{ijb} \leq L_b \, z_{ijb}$). La Etapa 2 recibe solamente el tiempo
de ruta de cada lote y decide qué picker hace cada lote y en qué orden, con el objetivo

\startformula
\min \left( W^{c} \cdot \sum_{b \in \mathcal{B}} c_b + W^{pen} \cdot \left( W^{e} \cdot \sum_{n \in \mathcal{N}} ea_n + W^{t} \cdot \sum_{n \in \mathcal{N}} ta_n \right) \right)
\stopformula

donde $c_b$ es el tiempo de completitud del lote, $ea_n$ el adelanto y $ta_n$ el atraso de cada orden.

La consecuencia central para la exploración es que las dos etapas tienen motores de costo
separados. La Etapa 1 no depende de cuántos pickers hay ni de sus velocidades: en el barrido de
pickers del estudio 01, su tiempo se mantuvo entre 13,9 s y 14,8 s con 3, 4 y 5 pickers. La
Etapa 2 depende del ruteo solo a través de los tiempos de ruta.
\stopsection

\startsection[title={Estudio 01: qué dimensión rompe primero}]
Desde una base de 6 lotes, 2 pickers, grilla 7×7, pesos 4/4 y velocidades uniformes, se movió un
eje por vez, con un tope de 300 s por resolución, fijando el mayor valor que resolvió antes de
pasar al eje siguiente.

\blank[medium]
\bTABLE[split=yes, align=middle, frame=on, width=fit]
\setupTABLE[row][1][background=color, backgroundcolor=darkgray, style=bold]
\bTR \bTH Eje \eTH \bTH Rango probado \eTH \bTH Resultado \eTH \eTR
\bTR \bTD Órdenes por lote \eTD \bTD 10 a 50 \eTD \bTD no rompió (a 50: Etapa 1 46,5 s, Etapa 2 14,1 s) \eTD \eTR
\bTR \bTD Lotes \eTD \bTD 6 a 8 \eTD \bTD rompe en 8: tope de 300 s con 518.628 nodos en la Etapa 2 \eTD \eTR
\bTR \bTD Pickers \eTD \bTD 2 a 5 \eTD \bTD no rompió, pero no es monótono \eTD \eTR
\bTR \bTD Grilla \eTD \bTD 7×7 a 15×15 \eTD \bTD rompe en 15×15: Etapa 1 de 464 s \eTD \eTR
\eTABLE
\blank[medium]

Rompen dos ejes, y en etapas distintas: los lotes rompen la Etapa 2 y la grilla rompe la
Etapa 1. Pasar de 6 a 8 lotes dio tope en los tres barridos donde se probó (pesos 4/4 con 26
órdenes por lote, pesos 2/6 con 50 y velocidades distintas con 50), así que es el eje más
restrictivo para la Etapa 2.

\startsubsection[title={Órdenes por lote: la Etapa 2 se vuelve más fácil}]
\placefigure[force][fig:ordenes_etapas]{Tiempo de cada etapa según las órdenes por lote (6 lotes, 2 pickers, grilla 7×7).}
{%%FIG:ordenes_etapas%%}

La Etapa 1 crece con las órdenes por lote, de 3,7 s a 46,5 s, porque cada lote tiene que pasar
por más ubicaciones. La Etapa 2 hace lo contrario de lo esperable: es más lenta con 10 órdenes
por lote (35,6 s) que con 26 (8,0 s). La figura \in[fig:nodos_sla] muestra que no es ruido de
reloj: el árbol de búsqueda se achica 37 veces, de 153.934 a 4.144 nodos.

\placefigure[force][fig:nodos_sla]{Nodos de la Etapa 2 según las órdenes por lote, con el porcentaje de órdenes a tiempo de cada punto.}
{%%FIG:nodos_sla%%}

La caída de nodos acompaña a la del SLA (de 96,7\% a 67,9\%). Una explicación posible es que,
con casi todas las órdenes a tiempo, los atrasos quedan en cero, el objetivo se aplana y el
solver no logra discriminar entre secuencias. Pero el SLA y las órdenes por lote están atados
(sección siguiente), así que el efecto no se puede atribuir a uno solo de los dos con estos
datos.
\stopsubsection

\startsubsection[title={Pickers: el caso difícil es el intermedio}]
\placefigure[force][fig:pickers]{Tiempo de la Etapa 2 según la cantidad de pickers, con 6 lotes en grilla 7×7.}
{%%FIG:pickers%%}

Con 6 lotes, el pico está en 4 pickers en los dos barridos, y los extremos son los más baratos:
2 pickers, o uno por lote. Una lectura posible: con pocos pickers hay mucho secuenciamiento pero
pocas asignaciones distintas, y con un picker por lote no hay nada que secuenciar; en el medio
se combinan las dos cosas.
\stopsubsection

\startsubsection[title={Grilla: el tamaño cuesta por sí mismo}]
Una grilla más grande hace que cada lote tenga más ubicaciones distintas, porque las órdenes
colisionan menos. Pero eso no explica el costo. Hay dos caminos para llegar a una cantidad
parecida de ubicaciones por lote, y cuestan órdenes de magnitud distintos:

\blank[medium]
\bTABLE[split=yes, align=middle, frame=on, width=fit]
\setupTABLE[row][1][background=color, backgroundcolor=darkgray, style=bold]
\bTR \bTH Camino \eTH \bTH Ubicaciones por lote \eTH \bTH Etapa 1 \eTH \eTR
\bTR \bTD 7×7, 26 órdenes por lote \eTD \bTD 18 a 22 \eTD \bTD 12,9 s \eTD \eTR
\bTR \bTD 7×7, 50 órdenes por lote \eTD \bTD 29 a 33 \eTD \bTD 46,5 s \eTD \eTR
\bTR \bTD 10×10, 26 órdenes por lote \eTD \bTD 21 a 26 \eTD \bTD 50,4 s \eTD \eTR
\bTR \bTD 15×15, 26 órdenes por lote \eTD \bTD 24 a 26 \eTD \bTD 464,4 s \eTD \eTR
\eTABLE
\blank[medium]

\placefigure[force][fig:dos_caminos]{Tiempo de la Etapa 1 contra ubicaciones por lote, según se llegue a ellas agrandando los pedidos o la grilla.}
{%%FIG:dos_caminos%%}

Llegar a unas 31 ubicaciones por lote agrandando los pedidos cuesta 46 s; llegar a 25
agrandando la grilla cuesta 464 s. El estudio 06 lo confirma desde el otro lado: en 25×25, con
solo 12 órdenes por lote, la Etapa 1 no terminó en 15 minutos. El mecanismo no está
identificado. Una sospecha sin verificar es que la relajación lineal del flujo, acotado por
$L_b$, se debilita a medida que crece el grafo.
\stopsubsection
\stopsection

\startsection[title={El SLA lo fija la función objetivo, no las fechas de entrega}]
Demorar un lote una unidad de tiempo cuesta $W^{c}$, más una unidad por cada orden suya que
ya está atrasada, y ahorra una unidad por cada orden que todavía llegaría antes de tiempo. En el
óptimo, el modelo demora cada lote mientras las órdenes adelantadas superen a las atrasadas en
más de $W^{c}$. Con $W^{e} = W^{t} = 1$ y fechas repartidas y alcanzables, eso da una
proporción de órdenes a tiempo de

\startformula
\mathrm{SLA} = \frac{N_b + W^{c}}{2 \, N_b}
\stopformula

donde $N_b$ es la cantidad de órdenes por lote.

\placefigure[force][fig:sla_formula]{SLA medido en el barrido de órdenes por lote contra el predicho por la fórmula.}
{%%FIG:sla_formula%%}

Los ocho puntos quedan a 1,7 puntos de la curva o menos, y el error se achica al crecer
$N_b$. El sesgo positivo viene de las órdenes que terminan exactamente en su fecha, que cuentan
como a tiempo. Tres consecuencias:

\startitemize[n]
\item {\bf Para diseñar escenarios no hace falta calibrar las fechas:} alcanza con que estén
repartidas y sean alcanzables. Con 26 órdenes por lote el SLA sale solo, cerca del 68\%.
\item {\bf El modelo es just-in-time:} agrega espera a propósito para no entregar antes de
tiempo. Eso lo vuelve un problema de scheduling con adelanto y atraso, que no es regular y es
una clase notoriamente difícil.
\item {\bf Los escenarios heredados de OPP\_v4 no cumplen la precondición:} todas las fechas son
30, antes de cualquier finalización posible, y por eso su SLA es 0\%.
\stopitemize
\stopsection

\startsection[title={Estudio 02: pesos de la grilla}]
Los escenarios heredados usan el mismo peso en todas las aristas (4/4), lo que produce muchísimas
rutas empatadas en costo. Se comparó contra pesos 2/6, elegidos para que la media siguiera siendo
4: lo único que cambia es la asimetría.

\placefigure[force][fig:pesos_ordenes]{Tiempo de la Etapa 1 con pesos uniformes y asimétricos, grilla 7×7.}
{%%FIG:pesos_ordenes%%}

En 7×7 los pesos asimétricos aceleran la Etapa 1 entre 1,7 y 3,7 veces, y la ventaja crece con
las órdenes por lote: sin empates, el solver no tiene que recorrer óptimos equivalentes para
probar que no hay nada mejor. La Etapa 2 sube y baja sin un patrón claro, salvo en el punto de
10 órdenes por lote, el de objetivo aplanado, donde baja de 35,6 s a 16,5 s.

\placefigure[force][fig:pesos_grilla]{Tiempo de la Etapa 1 según la grilla, con pesos uniformes y asimétricos (26 órdenes por lote).}
{%%FIG:pesos_grilla%%}

Pero la ventaja desaparece al crecer la grilla: en 15×15 los pesos 2/6 son incluso algo más
lentos (520 s contra 464 s). En una grilla grande los caminos mínimos ya toman muchos valores
distintos aunque los pesos sean uniformes, así que la asimetría no agrega nada. Los pesos
asimétricos ayudan donde no hacía falta y dejan de ayudar donde sí. Se adoptaron igual en el
perfil realista (4/5), por cobertura, porque un depósito real no tiene pasillos y travesaños del
mismo largo, pero no como optimización.
\stopsection

\startsection[title={Estudio 03: velocidades de los pickers}]
\startsubsection[title={Velocidades todas distintas}]
\placefigure[force][fig:velocidades_dispares]{Tiempo de la Etapa 2 con velocidades uniformes y con una velocidad distinta por picker.}
{%%FIG:velocidades_dispares%%}

Con 2 pickers, repartir las velocidades en $[0{,}7; 1{,}3]$ deja la Etapa 2 prácticamente
igual. La excepción otra vez es el punto de objetivo aplanado (10 órdenes por lote). Pero con 3
pickers y 50 órdenes por lote la Etapa 2 no terminó en 300 s (384.320 nodos), mientras que esa
misma forma con velocidades uniformes, en el barrido de pesos 2/6, había resuelto en 31 s.
\stopsubsection

\startsubsection[title={Velocidades en niveles}]
Una cuadrilla real no tiene una velocidad por persona: se contrata de un mismo grupo y todos
empujan el mismo carro, así que hay dos o tres bandas. Con la instancia fija (6 lotes, 26 órdenes
por lote, 7×7) se varió solo cuántas velocidades distintas hay, siempre dentro de
$[0{,}7; 1{,}3]$.

\placefigure[force][fig:niveles]{Tiempo de la Etapa 2 según cuántas velocidades distintas hay en la cuadrilla (escala logarítmica). \quote{tope}: no terminó en el límite; n/m: no medido.}
{%%FIG:niveles%%}

\blank[medium]
\bTABLE[split=yes, align=middle, frame=on, width=fit]
\setupTABLE[row][1][background=color, backgroundcolor=darkgray, style=bold]
\bTR \bTH Pickers \eTH \bTH Uniformes \eTH \bTH 2 niveles \eTH \bTH 3 niveles \eTH \bTH Todas distintas \eTH \eTR
\bTR \bTD 3 \eTD \bTD 17,7 s \eTD \bTD 134,2 s \eTD \bTD 113,5 s \eTD \bTD 112,7 s \eTD \eTR
\bTR \bTD 4 \eTD \bTD 26,3 s \eTD \bTD tope (549.456 nodos) \eTD \bTD tope (204.217) \eTD \bTD tope (347.556) \eTD \eTR
\bTR \bTD 6 \eTD \bTD 7,2 s \eTD \bTD 110,4 s \eTD \bTD 3.079 s, 499 nodos \eTD \bTD no medido \eTD \eTR
\eTABLE
\blank[medium]

\startitemize
\item {\bf El salto está entre una y dos velocidades.} Con 3 pickers, cualquier heterogeneidad
multiplica el tiempo por unas 7 veces; agregar niveles no suma.
\item {\bf Con 4 pickers ninguna configuración heterogénea termina en 300 s.} Es la
proporción intermedia que ya era la difícil en el estudio 01.
\item {\bf Con 6 pickers para 6 lotes} cada lote toma su picker y no hay secuenciamiento: la
heterogeneidad se tolera mejor, pero ese régimen es justamente el que no ejercita la Etapa 2.
\item {\bf Control de ruido:} con 3 pickers, \quote{3 niveles} y \quote{todas distintas} son la
misma configuración (0,7, 1,0 y 1,3) y dieron exactamente los mismos 250.037 nodos.
\stopitemize

Todos los bordes medidos antes, con velocidades uniformes, estaban sobre el caso
artificialmente fácil. Con velocidades realistas el alcance de la Etapa 2 es bastante menor.
Una explicación posible, sin verificar, es la simetría: pickers idénticos son intercambiables, y
eso le permite al solver descartar asignaciones equivalentes.
\stopsubsection
\stopsection

\startsection[title={Estudio 04: el tope duro funciona}]
Antes de calibrar se verificó que la herramienta corta cada punto a tiempo, en cualquier etapa.

\blank[medium]
\bTABLE[split=yes, align=middle, frame=on, width=fit]
\setupTABLE[row][1][background=color, backgroundcolor=darkgray, style=bold]
\bTR \bTH Punto \eTH \bTH Instancia \eTH \bTH Tope \eTH \bTH Resultado \eTH \eTR
\bTR \bTD T1, split \eTD \bTD 5×5, 4 lotes, 2 pickers \eTD \bTD 120 s \eTD \bTD óptimo en 3,2 s \eTD \eTR
\bTR \bTD T2, mono \eTD \bTD misma que T1 \eTD \bTD 120 s \eTD \bTD límite de SCIP a los 89,7 s, brecha 1,66\% \eTD \eTR
\bTR \bTD T3, split \eTD \bTD 7×7, 6 lotes, 4 pickers, 2 niveles \eTD \bTD 40 s \eTD \bTD límite de SCIP a los 27,7 s \eTD \eTR
\bTR \bTD T4, split \eTD \bTD 15×15, 6 lotes \eTD \bTD 45 s \eTD \bTD cortado en la Etapa 1 a los 45,3 s \eTD \eTR
\eTABLE
\blank[medium]

T4 es el caso que importa: la Etapa 1 no tiene límite propio, y el proceso se cortó a 0,3 s del
tope, dejando registrada la etapa en la que estaba.
\stopsection

\startsection[title={Estudio 05: calibración del perfil realista}]
El perfil que salió de los estudios anteriores: 6 lotes, 3 pickers, 26 órdenes por lote (156
órdenes), pesos 4/5 y dos bandas de velocidad (0,8 y 1,2). Solo varía el tamaño de la grilla.

\blank[medium]
\bTABLE[split=yes, align=middle, frame=on, width=fit]
\setupTABLE[row][1][background=color, backgroundcolor=darkgray, style=bold]
\bTR \bTH Grilla \eTH \bTH Ubic. por lote \eTH \bTH Etapa 1 \eTH \bTH Etapa 2 \eTH \bTH Total \eTH \bTH Nodos Etapa 2 \eTH \bTH SLA \eTH \eTR
\bTR \bTD 7×7 \eTD \bTD 18 a 22 \eTD \bTD 9,6 s \eTD \bTD 204,2 s \eTD \bTD 3,6 min \eTD \bTD 321.137 \eTD \bTD 67,9\% \eTD \eTR
\bTR \bTD 13×13 \eTD \bTD 23 a 26 \eTD \bTD 445,3 s \eTD \bTD 263,8 s \eTD \bTD 11,8 min \eTD \bTD 314.331 \eTD \bTD 67,9\% \eTD \eTR
\bTR \bTD 15×15 \eTD \bTD 24 a 26 \eTD \bTD más de 1.500 s \eTD \bTD -- \eTD \bTD cortado \eTD \bTD -- \eTD \bTD -- \eTD \eTR
\eTABLE
\blank[medium]

Medido de a un punto por vez, con la máquina sola. En 17×17 la Etapa 1 tampoco terminó en
25 minutos. Antes de fijar 6 lotes y 3 pickers se probaron otras formas de la Etapa 2 sobre 7×7.
Ninguna terminó en 25 minutos:

\placefigure[force][fig:formas]{Brecha de optimalidad al cabo de 25 minutos para distintas proporciones de lotes y pickers (7×7, perfil realista, en paralelo).}
{%%FIG:formas%%}

Con velocidades realistas, la Etapa 2 tolera unos dos lotes por picker; con más, o con la
proporción intermedia de 6 lotes para 4 pickers, ya no cierra. De esta calibración salieron los
tests de stress 04 a 07.
\stopsection

\startsection[title={Estudio 06: clausura métrica},reference=sec:clausura]
La clausura reducía el grafo de cada lote a sus ubicaciones requeridas, unidas por las
distancias de camino mínimo. Es exacta solo si la ruta, expandida de vuelta a la grilla, no
repite ningún vértice. Si repite, el lote se vuelve a resolver con el grafo completo.

En grillas con proporciones de depósito real (25×25 con 12 y 20 órdenes por lote, 35×35 con 12),
ninguna configuración terminó la Etapa 1 en 15 minutos, ni con clausura ni sin ella. Un
diagnóstico lote por lote mostró por qué:

\placefigure[force][fig:respaldo]{Lotes cuya ruta sobre la clausura repetía vértices y hubo que resolver de nuevo con el grafo completo.}
{%%FIG:respaldo%%}

Sobre la clausura, cada lote resolvía en décimas de segundo en 25×25 (0,16 s a 0,22 s). Pero en
el 67\% al 100\% de los lotes la ruta expandida repetía algún vértice, y el tiempo se iba en la
segunda resolución. En la calibración, con los mismos nodos en la Etapa 2, la Etapa 1 con
clausura fue más lenta que sin ella: 24,2 s contra 13,8 s en 7×7 y 737 s contra 676 s en 13×13.

La hipótesis, consistente con todo lo medido, es que en una grilla hay muchos caminos mínimos del
mismo largo: Dijkstra elige uno por tramo, dos tramos elegidos pueden cruzarse, y casi siempre
existía otro igual de corto que no se cruzaba. La clausura se eliminó del código.
\stopsection

\startsection[title={Estudio 07: la restricción de un arco por vértice}]
La Etapa 1 limita a uno los arcos que entran y salen de cada vértice
(\type{one_input_arc_per_node} y \type{one_output_arc_per_node}). Se probó quitarla, en una copia
temporal del modelo, para ver si cambiaba las respuestas o el tiempo.

\placefigure[force][fig:le1]{Tiempo de la Etapa 1 con y sin el límite de un arco por vértice (escala logarítmica).}
{%%FIG:le1%%}

\startitemize
\item {\bf Las respuestas no cambian.} Los tiempos de ruta de todos los lotes fueron idénticos
en las dos escalas, y los 11 tests funcionales que verifican la ruta exacta pasaron con el mismo
objetivo con y sin la restricción.
\item {\bf No es una palanca de rendimiento.} Sin ella, la Etapa 1 fue un 33\% más rápida en
7×7 y un 24\% más lenta en 13×13. En 25×25 no terminó en 15 minutos en ninguno de los dos casos.
\stopitemize
\stopsection

\startsection[title={Validación: los tests de stress del perfil realista},reference=sec:tests]
La corrida de la suite de stress del split del 29 de septiembre, con el código actual, la Etapa 1
en 6 procesos y $W^{c} = 27$ (mayor que las órdenes por lote, así que esperar a propósito no
conviene). La Etapa 1 se estima como el tiempo total menos el de la Etapa 2, así que incluye la
compilación y el resto de la preparación. La corrida anterior, del 23 de septiembre, usaba
$W^{c} = 9$: con los mismos tests, la Etapa 2 tardaba entre 86 y 203 s (salvo el test 11) y el
SLA quedaba en 68\%; con $W^{c} = 27$ la Etapa 2 baja a segundos y el SLA sube al 100\%.

\placefigure[force][fig:etapas_tests]{Tiempo de cada etapa en los tests de stress del perfil realista, con la Etapa 1 en 6 procesos.}
{%%FIG:etapas_tests%%}

\blank[medium]
\bTABLE[split=yes, align=middle, frame=on, width=fit]
\setupTABLE[row][1][background=color, backgroundcolor=darkgray, style=bold]
\bTR \bTH Test \eTH \bTH Grilla \eTH \bTH Etapa 1 (estim.) \eTH \bTH Etapa 2 \eTH \bTH Total \eTH \bTH SLA \eTH \eTR
\bTR \bTD 04 \eTD \bTD 7×7 \eTD \bTD 3,1 s \eTD \bTD 8,6 s \eTD \bTD 0,2 min \eTD \bTD 100\% \eTD \eTR
\bTR \bTD 05 \eTD \bTD 13×13 \eTD \bTD 95,1 s \eTD \bTD 9,7 s \eTD \bTD 1,7 min \eTD \bTD 100\% \eTD \eTR
\bTR \bTD 06 \eTD \bTD 14×14 \eTD \bTD 188,5 s \eTD \bTD 6,4 s \eTD \bTD 3,2 min \eTD \bTD 100\% \eTD \eTR
\bTR \bTD 07 \eTD \bTD 15×15 \eTD \bTD 883,6 s \eTD \bTD 4,2 s \eTD \bTD 14,8 min \eTD \bTD 100\% \eTD \eTR
\bTR \bTD 09 \eTD \bTD 8×21 \eTD \bTD 228,5 s \eTD \bTD 6,0 s \eTD \bTD 3,9 min \eTD \bTD 100\% \eTD \eTR
\bTR \bTD 11 \eTD \bTD 13×13, frío \eTD \bTD 107,0 s \eTD \bTD 7,2 s \eTD \bTD 1,9 min \eTD \bTD 97,4\% \eTD \eTR
\eTABLE
\blank[medium]

\placefigure[force][fig:etapa1_realista]{Tiempo de la Etapa 1 del perfil realista según la grilla (escala logarítmica).}
{%%FIG:etapa1_realista%%}

\startitemize
\item {\bf Con $W^{c} = 27$ domina la Etapa 1 en todos los tests}, aun en paralelo, y su
crecimiento es abrupto: de 13×13 a 14×14 se multiplica por 2,0; de 14×14 a 15×15, por 4,7.
\item {\bf La forma de la grilla pesa:} la grilla alargada de 8×21 (168 ubicaciones) tardó en
la Etapa 1 más que la cuadrada de 14×14 (196), aunque tiene menos ubicaciones. En la corrida del
23 de septiembre había quedado entre 13×13 y 14×14.
\item {\bf La Etapa 2 deja de pesar cuando esperar no conviene.} Con $W^{c} = 27$ resuelve en
menos de 10 s en todos los tests, también en el 11 (4 pickers, cadena de frío), donde ya era
barata con $W^{c} = 9$ porque el límite de frío deja elegibles solo a los 2 pickers rápidos.
\stopitemize
\stopsection

\startsection[title={La Etapa 1 en paralelo (estudio 12)}]
El split puede resolver en paralelo los ruteos de los lotes, que son independientes
(\type{OPP_STAGE1_WORKERS}; solo se activa con más de 4 lotes). Para medir cuánto aporta no
alcanza con comparar dos corridas de la suite, una en secuencia y otra en paralelo: entre corridas
distintas, la Etapa 2 de un mismo test (que no se paraleliza) llegó a variar el doble, tanto como
el efecto que se quiere medir.

Por eso se resolvió solo la Etapa 1 de las instancias exactas de los tests 05 (13×13) y 06
(14×14), con la máquina sola, en orden secuencial–paralelo–paralelo–secuencial para compensar el
calentamiento de la notebook.

\placefigure[force][fig:paralelo_ab]{Tiempo de la Etapa 1 en secuencia y con 6 procesos, dos corridas de cada una por instancia.}
{%%FIG:paralelo_ab%%}

\blank[medium]
\bTABLE[split=yes, align=middle, frame=on, width=fit]
\setupTABLE[row][1][background=color, backgroundcolor=darkgray, style=bold]
\bTR \bTH Instancia \eTH \bTH Secuencial \eTH \bTH 6 procesos \eTH \bTH Mejora \eTH \eTR
\bTR \bTD 13×13 (test 05) \eTD \bTD 241,8 s y 243,8 s \eTD \bTD 110,1 s y 110,8 s \eTD \bTD 2,2 veces \eTD \eTR
\bTR \bTD 14×14 (test 06) \eTD \bTD 401,4 s y 405,6 s \eTD \bTD 209,5 s y 215,9 s \eTD \bTD 1,9 veces \eTD \eTR
\eTABLE
\blank[medium]

Las dos corridas de cada modo quedaron a menos del 3\% entre sí, y las rutas fueron idénticas
en todas. {\bf El paralelismo divide por unas 2 el tiempo de la Etapa 1.} No llega a 6 por dos
razones: el tiempo en paralelo queda atado al lote más lento, y la notebook tiene 4 núcleos de
rendimiento y 8 de eficiencia, que van a cerca de la mitad de velocidad.

Sobre el test completo el efecto depende de cuánto pesa la Etapa 1: en 13×13 el test pasaría de
unos 6 minutos a unos 4, y en 15×15, donde la Etapa 1 es casi todo, de unos 28 a unos 15. En las
instancias chicas, donde manda la Etapa 2, no cambia nada. La memoria que reporta la
herramienta en paralelo (44 MB) es solo la del proceso principal: los procesos de cada lote no
se miden.
\stopsection

\startsection[title={Comparación puntual con el modelo monolítico},reference=sec:mono]
En dos puntos de la exploración se resolvió la misma instancia con OPP\_v5\_1\_mono.

\placefigure[force][fig:mono_split]{Split y mono sobre la misma instancia (escala logarítmica). El split terminó con el óptimo probado; el mono se detuvo por límite de tiempo.}
{%%FIG:mono_split%%}

\blank[medium]
\bTABLE[split=yes, align=middle, frame=on, width=fit]
\setupTABLE[row][1][background=color, backgroundcolor=darkgray, style=bold]
\bTR \bTH Instancia \eTH \bTH Split \eTH \bTH Mono \eTH \eTR
\bTR \bTD 5×5, 4 lotes, 2 pickers \eTD \bTD óptimo en 3,2 s (85 nodos) \eTD \bTD 89,7 s, brecha 1,66\% (25.867 nodos) \eTD \eTR
\bTR \bTD 7×7, perfil realista \eTD \bTD óptimo en 500 s en paralelo; 214 s solo \eTD \bTD 1.471 s, brecha 144,3\% \eTD \eTR
\eTABLE
\blank[medium]

En los dos casos el monolítico encontró una solución muy buena y no pudo demostrarlo. En 5×5 su
mejor solución ya era la óptima (8.966,27). En 7×7 quedó a menos del 1\% del óptimo (79.822,91
contra 79.046,86), pero su cota dual estaba cerca de 32.700, un 41\% del óptimo. La dificultad del
monolítico está en la cota, no en encontrar soluciones. Son dos puntos, corridos en condiciones
distintas; no alcanzan como comparación sistemática.
\stopsection

\startsection[title={Síntesis: cómo se comporta el modelo}]
\startitemize[n]
\item {\bf Las dos etapas se encarecen por motivos distintos.} La Etapa 1 depende del tamaño de
la grilla mucho más que de las ubicaciones por lote. La Etapa 2 depende de la proporción entre
lotes y pickers, de la heterogeneidad de velocidades y de qué tan plano queda el objetivo.
\item {\bf La Etapa 1 es lo que deja afuera a los depósitos reales.} Crece de forma abrupta con
la grilla: segundos en 7×7, minutos en 13×13, casi 15 minutos en 15×15 (en 6 procesos). En 25×25 no termina en
15 minutos ni con 12 órdenes por lote. El mecanismo no está identificado.
\item {\bf La Etapa 2 es lo que limita el tamaño de la cuadrilla.} Con velocidades realistas
tolera unos dos lotes por picker. La proporción intermedia (6 lotes para 4 pickers) es la más
difícil, y el salto de costo aparece apenas hay dos velocidades distintas.
\item {\bf El SLA no es un parámetro del escenario}, sino una consecuencia de $N_b$ y
$W^{c}$. Con $W^{c}$ menor que las órdenes por lote, el modelo retrasa lotes a propósito para
no adelantarse; con $W^{c} = 27$ eso deja de convenir y el SLA de la suite sube al 100\%.
\item {\bf Las simplificaciones de los escenarios heredados no son neutras:} los pesos uniformes
encarecen la Etapa 1 en grillas chicas; las velocidades uniformes abaratan la Etapa 2; las
fechas imposibles dejan la mitad del objetivo sin ejercitar.
\item {\bf Nada de lo probado para acelerar la Etapa 1 generalizó:} ni la clausura métrica, ni
los pesos asimétricos, ni quitar la restricción de un arco por vértice.
\item {\bf Frente al monolítico, el split prueba el óptimo donde el monolítico no cierra la
brecha}, en los dos puntos medidos.
\stopitemize
\stopsection

\startsection[title={Limitaciones y preguntas abiertas}]
\startitemize
\item Una sola semilla de generación: las conclusiones son sobre una familia de instancias, y
se vio que no se transfieren de una instancia a la de al lado sin volver a medir.
\item Los estudios 01 a 03 corrieron con la clausura activada y con uno o dos procesos en
paralelo; la calibración, con hasta ocho. Los tiempos absolutos de esos estudios no se deben
citar sin repetirlos.
\item El efecto del objetivo aplanado sobre la Etapa 2 estaba confundido con $N_b$. La suite
corrida con $W^{c} = 27$ sobre las mismas instancias ($N_b$ fijo) separa una parte: la Etapa 2
bajó de 86--203 s a menos de 10 s. Falta barrer valores intermedios de $W^{c}$.
\item Qué hace crecer la Etapa 1 con la grilla sigue sin explicación.
\item La franja de 10 a 20 minutos del perfil realista quedó cubierta con la corrida del 29 de
septiembre: 15×15 tardó 14,8 minutos (14×14, 3,2).
\stopitemize
\stopsection

\stoptext
